\section{Related work and claim boundary}

The ingredients of the workflow have substantial prior traditions: preregistration and prospective protocols, reproducible computational workflows, scientific provenance, negative-result reporting, autonomous research agents, iterative planning, verification pipelines, and donor/baseline comparison. Q2 claims none of these ingredients in isolation.

The candidate residual is their executable composition around recursive negative-result recovery: a typed terminal system in which donor absorption, refutation, support, and cannot-check remain distinct; reopening is tied to a recorded residual; and the evidence path is receipt-bound. That residual is conditional on a fresh submission-date literature search. The manuscript therefore avoids `first`, `unique`, or exhaustive-nearest-work language until that search is complete.

Q1 owns the exact TARE expressivity result. Q3 owns the controller--host agreement benchmark and is now bounded away from the systems-paper territory being opened by P15. Q4 owns the typed/scoped partial-knowledge mechanism evidence. Q2 cites those outcomes only to show the behavior of the common methodology.